% \begin{name}
% 	{GV biên soạn: Huỳnh Văn Quy\\ GV phản biện: Thầy Trí Thiện}
% 	{TOÁN 10-KNTT-CS}
% \end{name}
%----------------NỘI DUNG BIÊN SOẠN
% \setcounter{chapter}{0}
\chap{MỆNH ĐỀ - TẬP HỢP}
\setcounter{section}{0}
\def\tenchude{MỆNH ĐỀ}
\section{MỆNH ĐỀ}

\subsection{KIẾN THỨC TRỌNG TÂM}
\subsubsection{Mệnh đề}
\begin{khung4}{}
	\begin{itemize}
		\item Mệnh đề là một khẳng định hoặc \textbf{đúng} hoặc \textbf{sai}.
		\item Một mệnh đề không thể vừa đúng vừa sai.
	\end{itemize}
\end{khung4}
\begin{note}
	 Người ta thường sử dụng các chữ cái in hoa $P$, $Q$, $R$, \ldots  để kí hiệu mệnh đề.
\end{note}
\subsubsection{Mệnh đề chứa biến}
\begin{khung4}{}
	Những khẳng định mà tính đúng, sai của chúng phụ thuộc vào giá trị của biến gọi là \textbf{mệnh đề chứa biến}.
\end{khung4}
\subsubsection{Mệnh đề phủ định}
\begin{khung4}{}
	Mỗi mệnh đề $P$ có \textbf{mệnh đề phủ định}, kí hiệu là $\overline{P}$.\\
	Mệnh đề $P$ và mệnh đề phủ định $\overline{P}$ của nó có tính đúng sai trái ngược nhau. \\Nghĩa là khi $P$ đúng thì $\overline{P}$ sai, khi $P$ sai thì $\overline{P}$ đúng.
\end{khung4}
\subsubsection{Mệnh đề kéo theo}
\begin{khung4}{}
	Cho hai mệnh đề $P$ và $Q$. Mệnh đề \lq\lq Nếu $P$ thì $Q$\rq\rq\ được gọi là mệnh đề kéo theo.\\
	Mệnh đề $P\Rightarrow Q$ chỉ sai khi $P$ đúng và $Q$ sai.
\end{khung4}
\begin{nx}
	 Mệnh đề $P\Rightarrow Q$ còn được phát biểu là \lq\lq $P$ kéo theo $Q$\rq\rq, \lq\lq $P$ suy ra $Q$\rq\rq.
\end{nx}
\begin{note}
	Trong toán học, định lí là một mệnh đề đúng, thường có dạng $P\Rightarrow Q$.
	Khi đó ta nói
	\begin{enumerate}[•]
		\item $P$ là giả thiết, $Q$ là kết luận của định lí.
		\item $P$ là \textbf{điều kiện đủ}  để có $Q$, còn $Q$ là  \textbf{điều kiện cần}  để có $P$.
	\end{enumerate}
\end{note}
\subsubsection{Mệnh đề đảo. Hai mệnh đề tương đương}
\begin{khung4}{}
	Mệnh đề $Q\Rightarrow P$ được gọi là \textbf{mệnh đề đảo} của mệnh đề $P\Rightarrow Q$.
\end{khung4}
\begin{note}
	Mệnh đề đảo của một mệnh đề đúng không nhất thiết là đúng.
\end{note}
\begin{khung4}{}
	Nếu cả hai mệnh đề $P \Rightarrow Q$ và $Q \Rightarrow P$ đều đúng thì ta nói $P$ và $Q$ là hai \textbf{mệnh đề tương đương}, kí hiệu là $P \Leftrightarrow Q$ (đọc là \lq\lq $P$ tương đương $Q$\rq\rq~hoặc \lq\lq $P$ khi và chỉ khi $Q$\rq\rq).\\
	Khi đó, ta cũng nói $P$ là \textbf{điều kiện cần và đủ} để có $Q$ (hay $Q$ là điều kiện cần và đủ để có $P$).
\end{khung4}
\begin{nx}
	Hai mệnh đề $P$ và $Q$ tương đương khi chúng cùng đúng hoặc cùng sai.
\end{nx}

\subsubsection{Mệnh đề chứa kí hiệu $\forall$ và $\exists$}
\begin{khung4}{}
	Mệnh đề \lq\lq $\forall x\in M, P(x)$\rq\rq~ đúng nếu với mọi $x_0\in M$, $P\left(x_0\right)$ là mệnh đề đúng.\\
	Mệnh đề \lq\lq $\exists x\in M, P(x)$\rq\rq~ đúng nếu có $x_0\in M$ sao cho $P\left(x_0\right)$ là mệnh đề đúng.
\end{khung4}
%-------------------------------------------------------------------------------------------------------------
\subsection{PHÂN LOẠI VÀ PHƯƠNG PHÁP GIẢI TOÁN}
\begin{dang}{Xác định mệnh đề phủ định và xét tính đúng sai của một mệnh đề}
	 \textbf{Phương pháp giải:}
	\begin{enumerate}
		\item \textbf{Bước 1: Tìm mệnh đề phủ định.}
		\begin{itemize}
			\item Để phủ định một mệnh đề, ta thường thêm (hoặc bớt) từ \lq\lq không\rq\rq\ hoặc \lq\lq không phải\rq\rq\ vào trước vị ngữ.
			\item Phủ định của mệnh đề \lq\lq $\forall x, P(x)$ \rq\rq\ là \lq\lq $\exists x, \text{ không } P(x)$\rq\rq.
			\item Phủ định của mệnh đề \lq\lq $\exists x, P(x)$ \rq\rq\ là \lq\lq $\forall x, \text{ không } P(x)$\rq\rq.
		\end{itemize}
		\item \textbf{Bước 2: Xét tính đúng sai của mệnh đề.}
		\begin{itemize}
			\item Một mệnh đề và mệnh đề phủ định của nó có tính đúng sai trái ngược nhau.
			\item Nếu mệnh đề ban đầu đúng thì mệnh đề phủ định sai và ngược lại.
		\end{itemize}
	\end{enumerate}
\end{dang}

 \begin{vd}%[Thành Đức Trung]%[0D1Y1-2]
	Đánh dấu \lq\lq  x\rq\rq\ vào ô thích hợp trong bảng sau
	\begin{center}
		\begin{tabular}{|>{\centering\arraybackslash}m{5.5cm}|>{\centering\arraybackslash}m{2cm}|>{\centering\arraybackslash}m{2cm}|>{\centering\arraybackslash}m{2cm}|}
			\hline
			Câu & Không phải MĐ & MĐ đúng & MĐ sai \\
			\hline
		\end{tabular}
		\begin{tabular}{|>{\raggedright\arraybackslash}m{5.5cm}|>{\centering\arraybackslash}m{2cm}|>{\centering\arraybackslash}m{2cm}|>{\centering\arraybackslash}m{2cm}|}
			$13$ là số nguyên tố. &  &  &  \\
			\hline
			Tổng độ dài hai cạnh bất kì của một tam giác nhỏ hơn độ dài cạnh còn lại. &  &  &  \\
			\hline
			Bạn đã làm bài tập chưa? &  &  &  \\
			\hline
			Thời tiết hôm nay thật đẹp! &  &  &  \\
			\hline
			$9>2$. &  &  &  \\
			\hline
			$27$ chia hết cho $5$. &  &  &  \\
			\hline
			$2+3=6$. &  &  &  \\
			\hline
			$36$ là số chính phương. &  &  &  \\
			\hline
			Chó có khôn hơn lợn không? &  &  &  \\
			\hline
		\end{tabular}
	\end{center}
	\loigiai
	{
		\begin{center}
			\begin{tabular}{|>{\centering\arraybackslash}m{5.5cm}|>{\centering\arraybackslash}m{2cm}|>{\centering\arraybackslash}m{2cm}|>{\centering\arraybackslash}m{2cm}|}
				\hline
				Câu & Không phải mệnh đề & Mệnh đề đúng & Mệnh đề sai \\
				\hline
			\end{tabular}
			\begin{tabular}{|>{\raggedright\arraybackslash}m{5.5cm}|>{\centering\arraybackslash}m{2cm}|>{\centering\arraybackslash}m{2cm}|>{\centering\arraybackslash}m{2cm}|}
				$13$ là số nguyên tố. &  & x &  \\
				\hline
				Tổng độ dài hai cạnh bất kì của một tam giác nhỏ hơn độ dài cạnh còn lại. &  &  & x \\
				\hline
				Bạn đã làm bài tập chưa? & x &  &  \\
				\hline
				Thời tiết hôm nay thật đẹp! & x &  &  \\
				\hline
				$9>2$. &  & x &  \\
				\hline
				$27$ chia hết cho $5$. &  &  & x \\
				\hline
				$2+3=6$. &  &  & x \\
				\hline
				$36$ là số chính phương. &  & x &  \\
				\hline
				Chó có khôn hơn lợn không? & x &  &  \\
				\hline
			\end{tabular}
		\end{center}
	}
\end{vd}
%%%=============VD_1=============%%%
\begin{vd}%[0D1H1-3]%
	Phát biểu mệnh đề phủ định của mỗi mệnh đề sau:	
		\begin{enumerate}
			\item $25$ là số chính phương \dotfill;
			\item Hình chữ nhật không phải là hình vuông\dotfill.
		\end{enumerate}
\loigiai{
		\begin{enumerate}
			\item $25$ không phải là số chính phương;
			\item Hình chữ nhật là hình vuông.
		\end{enumerate}
	}
\end{vd}
\begin{vd}%[0D1H1-2]
	Xác định tính đúng sai của mỗi mệnh đề sau
	\begin{enumerate}
		\item $\pi<\dfrac{10}{3}$\dotfill.
		\item Phương trình $3x+7=0$ có nghiệm\dotfill.
		\item Có ít nhất một số cộng với chính nó bằng $0$\dotfill.
		\item $2022$ là hợp số\dotfill.
	\end{enumerate}
\loigiai{
	\begin{enumerate}
		\item \lq\lq $\pi<\dfrac{10}{3}$\rq\rq\, là mệnh đề đúng vì $\pi\approx 3{,}14 <3{,}33\approx \dfrac{10}{3}$.
		\item \lq\lq Phương trình $3x+7=0$ có nghiệm\rq\rq\, là mệnh đề đúng vì $3x+7=0 \Leftrightarrow x=-\dfrac{7}{3}$.
		\item \lq\lq Có ít nhất một số cộng với chính nó bằng $0$\rq\rq\, là mệnh đề đúng vì $0+0=0$.
		\item \lq\lq $2022$ là hợp số\rq\rq\, là mệnh đề đúng vì $2022=2\cdot3\cdot337$.
	\end{enumerate}
	}
\end{vd}
\begin{vd}%[0D1H1-2]%
	Xét tính đúng sai của các mệnh đề sau và phát biểu mệnh đề phủ định của chúng.
	\begin{enumerate}[a)]
		\item \lq\lq$2019$ chia hết cho $3$\rq\rq;\dongcham{2}
		\item \lq\lq$\pi<3{,}15$\rq\rq;\dongcham{2}
		\item \lq\lq Tam giác có hai góc bằng $45^{\circ}$ là tam giác vuông cân\rq\rq.\dongcham{2}
	\end{enumerate}
\loigiai{
	\begin{enumerate}[a)]
		\item \lq\lq$2019$ chia hết cho $3$\rq\rq\ là mệnh đề đúng.\\
		Mệnh đề phủ định là \lq\lq$2019$ không chia hết cho $3$\rq\rq.
		\item \lq\lq$\pi<3{,}15$\rq\rq\ là mệnh đề đúng.\\
		Mệnh đề phủ định là \lq\lq $\pi \geq 3{,}15$\rq\rq.
		\item \lq\lq Tam giác có hai góc bằng $45^{\circ}$ là tam giác vuông cân\rq\rq\ là mệnh đề đúng.\\
		Mệnh đề phủ định là \lq\lq Tam giác có hai góc bằng $45^{\circ}$ không là tam giác vuông cân\rq\rq.
	\end{enumerate}
	}
\end{vd}
\begin{vd}%[0D1H1-2]
		Cho biểu thức $P(n)=3n^2+3n-6$ với $n\in \mathbb{N}$. Xét tính đúng - sai của các phát biểu sau:
		\choiceTF
		{$P(3)$ có giá trị là số lẻ}
		{$P(4)+1$ có giá trị là số nguyên tố}
		{\True $P(7n)<P(n)+2025$ với $n=3$}
		{Tồn tại số tự nhiên $n$ để biểu thức $\dfrac{P(n)+1}{n+2}$ có giá trị nguyên}
		\loigiai{
			\begin{itemchoice}
				\itemch Vì $P(3)=30$ là giá trị chẵn.
				\itemch Vì $P(4)+1=55$ không phải là số nguyên tố
				\itemch Vì với $n=3\Rightarrow P(7n)=P(21)=1380$; $P(n)+2025=P(3)+2025=2055$.\\ Suy ra $P(7n)<P(n)+2025$.
				\itemch Ta có $\dfrac{P(n)+1}{n+2}=\dfrac{3n^2+3n-5}{n+2}=3n-3+\dfrac{1}{n+2}$. \\
				Khi đó $\dfrac{P(n)+1}{n+2} \in \mathbb{Z} \Leftrightarrow \dfrac{1}{n+1} \in \mathbb{Z} \Leftrightarrow (n+1) \in \{-1;1\} \Leftrightarrow n \in \{-3; -1\} \not \subset \mathbb{N}$.
			\end{itemchoice}
		}
	\end{vd}
\begin{dang}{Xác định mệnh đề kéo theo, mệnh đề đảo, mệnh đề tương đương}
	\textbf{Phương pháp giải:}
	Cho hai mệnh đề $P$ và $Q$.
	\begin{itemize}
		\item \textbf{Mệnh đề kéo theo:} có dạng \lq\lq Nếu $P$ thì $Q$\rq\rq\, ký hiệu là $P \Rightarrow Q$.
		\begin{itemize}
			\item Mệnh đề này chỉ sai khi $P$ đúng và $Q$ sai.
		\end{itemize}
		\item \textbf{Mệnh đề đảo:} là mệnh đề \lq\lq Nếu $Q$ thì $P$\rq\rq ký hiệu là $Q \Rightarrow P$.
		\begin{itemize}
			\item Lưu ý: Tính đúng sai của mệnh đề đảo ($Q \Rightarrow P$) không phụ thuộc vào tính đúng sai của mệnh đề ban đầu ($P \Rightarrow Q$).
		\end{itemize}
		\item \textbf{Mệnh đề tương đương:} có dạng \lq\lq $P$ nếu và chỉ nếu $Q$\rq\rq, ký hiệu là $P \Leftrightarrow Q$.
		\begin{itemize}
			\item Mệnh đề này đúng khi cả hai mệnh đề kéo theo thuận ($P \Rightarrow Q$) và đảo ($Q \Rightarrow P$) đều đúng.
		\end{itemize}
	\end{itemize}
\end{dang}
\setcounter{vd}{0}
%%%=============VD_1=============%%%
\begin{vd}%[0D1N1-4]
	Cho tam giác $ABC$. Xét hai mệnh đề:\\
	$P$: \lq \lq Tam giác $ABC$ có hai góc bằng $60^\circ$\rq \rq .\\
	$Q$: \lq \lq Tam giác $ABC$ đều\rq \rq. \\
	Hãy phát biểu mệnh đề $P \Rightarrow Q$ và nhận xét tính đúng sai của mệnh đề đó.\\
	\dongcham{2}
\loigiai{
	$P \Rightarrow Q$: \lq \lq Nếu tam giác $ABC$ có hai góc bằng $60^\circ$ thì tam giác $ABC$ đều\rq \rq.\\
	Mệnh đề kéo theo này là mệnh đề đúng.
	}
\end{vd}
%%%=============================%%%
%%%=============VD_2=============%%%
\begin{vd}%[0D1N1-4]%
	Cho $n$ là số tự nhiên. Xét các mệnh đề\\
	$P\colon$ \lq\lq $n$ là một số tự nhiên chia hết cho $16$\rq\rq,\\
	$Q\colon$ \lq\lq $n$ là một số tự nhiên chia hết cho $8$\rq\rq.
	\begin{enumerate}[a)]
		\item Phát biểu mệnh đề $P\Rightarrow Q$. Nhận xét tính đúng sai của mệnh đề đó.\dongcham{3}
		\item Phát biểu mệnh đề đảo của mệnh đề $P\Rightarrow Q$. Nhận xét tính đúng sai của mệnh đề đó.\dongcham{3}
	\end{enumerate}
\loigiai{
	\begin{enumerate}[a)]
		\item Mệnh đề $P\Rightarrow Q\colon$ \lq\lq Nếu số tự nhiên $n$ chia hết cho $16$ thì $n$ chia hết cho $8$\rq\rq. Đây là mệnh đề đúng vì $8$ là ước của $16$.
		\item Mệnh đề đảo của mệnh đề $P\Rightarrow Q$ là mệnh đề $Q\Rightarrow P\colon$ \lq\lq Nếu số tự nhiên $n$ chia hết cho $8$ thì $n$ chia hết cho $16$\rq\rq. Đây là mệnh đề sai vì với $n=8$, $n$ chia hết cho $8$ nhưng không chia hết cho $16$.
	\end{enumerate}
	}
\end{vd}
%%%=============================%%%
%%%=============VD_3=============%%%
\begin{vd}%[0D1H1-4]
	Xét hai mệnh đề:\\
	$P\colon$ \lq\lq Tứ giác $A B C D$ là hình bình hành\rq\rq;\\
	$Q\colon$ \lq\lq Tứ giác $A B C D$ có hai đường chéo cắt nhau tại trung điểm của mỗi đường\rq\rq.
	\begin{enumerate}[a)]
		\item Phát biểu mệnh đề $P \Rightarrow Q$ và xét tính đúng sai của nó.\dongcham{0}
		\item Phát biểu mệnh đề đảo của mệnh đề $P \Rightarrow Q$.\dongcham{2}
	\end{enumerate}
\loigiai{
	\begin{enumerate}[a)]
		\item $P \Rightarrow Q\colon$ \lq\lq Nếu tứ giác $ABCD$ là hình bình hành thì tứ giác $ABCD$ có hai đường chéo cắt nhau tại trung điểm mỗi đường\rq\rq.\\
		Mệnh đề $P \Rightarrow Q$ đúng.
		\item $Q \Rightarrow P\colon$ \lq\lq Nếu tứ giác $ABCD$ có hai đường chéo cắt nhau tại trung điểm của mỗi đường thì tứ giác $ABCD$ là hình bình hành\rq\rq.
	\end{enumerate}
	}
\end{vd}
%%%=============================%%%
\begin{dang}{Xác định mệnh đề phủ định của mệnh đề chứa kí hiệu $\forall$ và $\exists$}
	\begin{itemize}
		\item Ta đọc $\forall x\in X$, $\exists x\in X$ lần lượt là với mọi $x$ thuộc $X$, tồn tại (có ít nhất) $x$ thuộc $X$.
		\item Phủ định của mệnh đề \lq\lq$\forall x\in X,P(x )$\rq\rq\,là mệnh đề \lq\lq$\exists x\in X,\overline{P(x)}$\rq\rq.
		\item Phủ định của mệnh đề \lq\lq$\exists x\in X,P(x )$\rq\rq\,là mệnh đề \lq\lq$\forall x\in X,\overline{P(x)}$\rq\rq.
	\end{itemize}
\end{dang}
\setcounter{vd}{0}
%%%=============VD_1=============%%%
\begin{vd}%[0D1H1-5]
	Lập mệnh đề phủ định của mỗi mệnh đề sau và xét tính đúng sai của mỗi mệnh đề phủ định đó:
		\begin{enumerate}[a)]
			\item $\forall x \in \mathbb{R},~ x^2 \neq 2x-2$ \dotfill.
			\item $\forall x \in \mathbb{R},~ x^2\leq 2x-1$ \dotfill.
			\item $\exists x \in \mathbb{R},~ x+\dfrac{1}{x} \geq 2$ \dotfill.
			\item $\exists x \in \mathbb{R},~ x^2-x+1< 0$ \dotfill.
		\end{enumerate}
\loigiai{
	\begin{enumerate}[a)]
		\item Phủ định của mệnh đề \lq\lq $\forall x \in \mathbb{R}, x^2\neq 2x-2$\rq\rq~là mệnh đề \lq\lq $\exists x \in \mathbb{R}, x^2=2x-2$\rq\rq. \\
		Mệnh đề phủ định sai vì phương trình $x^2=2x-2$ vô nghiệm nên không có giá trị nào của $x$ thoả mãn $x^2=2x-2$.
		\item Phủ định của mệnh đề \lq\lq $\forall x \in \mathbb{R}, x^2\leq 2x-1$\rq\rq~là mệnh đề \lq\lq $\exists x \in \mathbb{R}, x^2> 2x-1$ \rq\rq. \\
		Mệnh đề phủ định đúng vì với $x=2$, ta có $2^2> 2\cdot 2-1$.
		\item Phủ định của mệnh đề \lq\lq $\exists x \in \mathbb{R}, x+\dfrac{1}{x} \geq 2$\rq\rq~là mệnh đề \lq\lq $\forall x \in \mathbb{R}, x+\dfrac{1}{x} < 2$ \rq\rq.\\
		Mệnh đề phủ định sai vì với $x=2$, ta có $2+\dfrac{1}{2} > 2$.
		\item Phủ định của mệnh đề \lq\lq $\exists x \in \mathbb{R}, x^2-x+1< 0$\rq\rq~là mệnh đề \lq\lq $\forall x \in \mathbb{R}, x^2-x+1\geq 0$\rq\rq.\\
		Mệnh đề phủ định đúng vì $x^2-x+1=\left(x-\dfrac{1}{2}\right)^2+\dfrac{3}{4} > 0, \forall x \in \mathbb{R}$.
	\end{enumerate}
	}
\end{vd}
%%%=============================%%%
%%%=============VD_2=============%%%
\begin{vd}%[0D1H1-5]
	Cho các mệnh đề sau:\\
	$P\colon$ \lq\lq Giá trị tuyệt đối của mọi số thực đều lớn hơn hoặc bằng chính nó\rq\rq;\\
	$Q\colon$ \lq\lq Có số tự nhiên sao cho bình phương của nó bằng $10$ \rq\rq;\\
	$R\colon$ \lq\lq Có số thực $x$ sao cho $x^2+2x-1=0$\rq\rq.
	\begin{enumerate}[a)]
		\item Xét tính đúng sai của mỗi mệnh đề trên.
		\item Sử dụng kí hiệu $\forall$, $\exists$ để viết lại các mệnh đề đã cho.
	\end{enumerate}
\dongcham{4}
\loigiai{
	\begin{enumerate}[a)]
		\item $P\colon$ Đúng vì với $x \ge 0$ thì $|x|=x$ và với $x<0$ thì $|x|>0>x$.\\
		$Q\colon$ Sai vì $x^2=10 \Leftrightarrow x=\pm \sqrt{10}$, mà $\sqrt{10}\notin \mathbb{N}$.\\
		$R\colon$ Đúng vì phương trình $x^2+2x-1=0$ có nghiệm $x=-1+\sqrt{2}$ và $x=-1-\sqrt{2}$.
		\item $P\colon$ \lq\lq$\forall x\in\mathbb{R}, |x|\geq x$\rq\rq;\\ 
		$Q\colon$ \lq\lq$\exists x\in\mathbb{N}, x^2=10$\rq\rq;\\
		$R\colon$ \lq\lq$\exists x\in\mathbb{R}, x^2+2x-1=0$\rq\rq.
	\end{enumerate}
	}
\end{vd}
%%%=============================%%%
%%%=============VD_3=============%%%
\begin{vd}%[0D1H1-5]
	Xét tính đúng sai và viết mệnh đề phủ định của các mệnh đề sau:
		\begin{enumerate}[a)]
			\item \lq\lq$\forall x\in\mathbb{R}, x^2+2x+2>0$\rq\rq\dongcham{2}.
			\item \lq\lq$\exists x\in\mathbb{R}, x^2+3x+4=0$\rq\rq\dongcham{2}.
		\end{enumerate}
\loigiai{
	\begin{enumerate}[a)]
		\item Mệnh đề đúng, vì $x^2+2x+2=\left(x^2+2x+1\right)+1=(x+1)^2+1>0$ với mọi số thực $x$.\\
		Mệnh đề phủ định của mệnh đề này là \lq\lq$\exists x\in\mathbb{R}, x^2+2x+2\leq 0$\rq\rq.
		\item Mệnh đề sai, vì phương trình $x^2+3x+4=0$ vô nghiệm $(\Delta=-7<0)$.\\
		Mệnh đề phủ định của mệnh đề này là \lq\lq$\forall x\in\mathbb{R}, x^2+3x+4\neq 0$\rq\rq.
	\end{enumerate}
	}
\end{vd}
%%%=============================%%%
\subsection{BÀI TẬP TỰ LUYỆN}
\baitaptl
\setcounter{bt}{0}
%%%=============EX_1=============%%%
\begin{bt}%[0D1N1-3]
	Phát biểu mệnh đề phủ định của các mệnh đề sau:
	\begin{enumerate}[a)]
		\item $P\colon$ \lq\lq Tháng $12$ dương lịch có $30$ ngày\rq\rq;
		\item $Q\colon$ \lq\lq $9^{10}\ge 10^9$\rq\rq;
		\item $R\colon$ \lq\lq Phương trình $x^2+1=0$ có nghiệm\rq\rq.
	\end{enumerate}
\loigiai{Mệnh đề phủ định của các mệnh đề trên là
	\begin{enumerate}[a)]
		\item $\overline P\colon$ \lq\lq Tháng $12$ dương lịch không phải có $30$ ngày\rq\rq;
		\item $\overline Q\colon$ \lq\lq $9^{10}<10^9$\rq\rq;
		\item $\overline R\colon$ \lq\lq Phương trình $x^2+1=0$ vô nghiệm\rq\rq.
	\end{enumerate}
	}
\end{bt}
\begin{bt}%[0D1H1-5]
	Xét tính đúng sai và viết mệnh đề phủ định của các mệnh đề sau đây:
	\begin{multicols}{3}
		\begin{enumerate}[a)]
			\item \lq\lq$\exists x\in\mathbb{N}, x+3=0$\rq\rq;
			\item \lq\lq$\forall x\in\mathbb{R}, x^2+1\geq 2x$\rq\rq;
			\item \lq\lq$\forall a\in\mathbb{R},\sqrt{a^2}=a$\rq\rq.
		\end{enumerate}
	\end{multicols}
\loigiai{
	\begin{enumerate}[a)]
		\item \lq\lq$\exists x\in\mathbb{N}, x+3=0$\rq\rq ~là mệnh đề sai. Mệnh đề phủ định là \lq\lq$\forall x\in\mathbb{N}, x+3\neq 0$\rq\rq.
		\item \lq\lq$\forall x\in\mathbb{R}, x^2+1\geq 2x$\rq\rq~là mệnh đề đúng. Mệnh đề phủ định là \lq\lq$\exists x\in\mathbb{R}, x^2+1<2x$\rq\rq.
		\item \lq\lq$\forall a\in\mathbb{R},\sqrt{a^2}=a$\rq\rq~là mệnh đề sai. Mệnh đề phủ định là \lq\lq$\exists a\in\mathbb{R},\sqrt{a^2}\neq a$\rq\rq.
	\end{enumerate}
	}
\end{bt}
\begin{bt}%[0D1H1-4]
	Cho các định lí:\\
	$P\colon$ \lq\lq Nếu hai tam giác bằng nhau thì diện tích của chúng bằng nhau\rq\rq;\\
	$Q\colon$ \lq\lq Nếu $a<b$ thì $a+c<b+c$ \rq\rq~($a$, $b$, $c \in \mathbb{R}$).
	\begin{enumerate}[a)]
		\item Chỉ ra giả thiết và kết luận của mỗi định lí.
		\item Phát biểu lại mỗi định lí đã cho, sử dụng thuật ngữ \lq\lq điều kiện cần\rq\rq~ hoặc \lq\lq ~điều kiện đủ\rq\rq.
		\item Mệnh đề đảo của mỗi định lí đó có là định lí không?
	\end{enumerate}
\loigiai{
	\begin{enumerate}[a)]
		\item Mệnh đề $P$.\\
		Giả thiết: hai tam giác bằng nhau.\\
		Kết luận: diện tích của chúng bằng nhau.\\
		Mệnh đề $Q$.\\
		Giả thiết: $a<b$.\\
		Kết luận: $a+c<b+c$ ($a$, $b$, $c \in \mathbb{R}$).
		\item Phát biểu định lí:
		\begin{itemize}
			\item $P\colon$\lq\lq Hai tam giác bằng nhau là điều kiện đủ để diện tích của chúng bằng nhau\rq\rq~hoặc~\lq\lq Diện tích của hai tam giác bằng nhau là điều kiện cần để hai tam giác bằng nhau\rq\rq.
			\item $Q\colon$ \lq\lq $a<b$ là điều kiện đủ để $a+c<b+c$ ($a$, $b$, $c \in \mathbb{R}$)\rq\rq~hoặc~\lq\lq $a+c<b+c$ ($a$, $b$, $c \in \mathbb{R}$)\rq\rq~là điều kiện cần để \lq\lq$a<b$\rq\rq.
		\end{itemize}
		\item Mệnh đề đảo của mỗi định lí là mệnh đề sai nên không là định lí.
	\end{enumerate}
	}
\end{bt}
\begin{bt}%[0D1H1-2]
	Cho các mệnh đề chứa biến:
	\begin{enumerate}[a)]
		\item $P(x)\colon$\lq\lq $2x=1$\rq\rq;
		\item $R(x,y)\colon$\lq\lq $2x+y=3 $ \rq\rq~(mệnh đề này chứa hai biến $x$ và $y$);
		\item $T(n)\colon$\lq\lq $2n+1$ là số chẵn\rq\rq~($n$ là số tự nhiên).
	\end{enumerate}
	Với mỗi mệnh đề chứa biến trên, tìm những giá trị của biến để nhận được một mệnh đề đúng và một mệnh đề sai.
\loigiai{
	\begin{enumerate}[a)]
		\item Với $x=\dfrac{1}{2}$ thì $P\left(\dfrac{1}{2}\right)\colon$
		\lq\lq$2\cdot\dfrac{1}{2}=1$\rq\rq~là mệnh đề đúng.\\
		Với $x=1$ thì $P(1)\colon$\lq\lq $2\cdot1=1$\rq\rq~là mệnh đề sai.
		\item Với $x=1$, $y=1$ thì $R(1,1)\colon$\lq\lq $2\cdot1+1=3$\rq\rq~là mệnh đề đúng.\\
		Với $x=1$, $y=1$ thì $R(1,2)\colon$\lq\lq $2\cdot1+2=3$\rq\rq~là mệnh đề sai.
		\item Lấy số tự nhiên $n_0$ bất kì ta đều được $2n_0+1$ là một số lẻ, nghĩa là $T\left(n_0\right)\colon$\lq\lq $2n_0+1$ là số chẵn \rq\rq~ là mệnh đề sai. Do đó, không có giá trị n$_0$ của $n$ để $T\left(n_0\right)$ là mệnh đề đúng. $T\left(n_0\right)$ là mệnh đề sai với số tự nhiên $n_0 $ bất kì.
	\end{enumerate}
	}
\end{bt}
\begin{bt}%[0D1H1-4]
	Sử dụng các thuật ngữ \lq\lq điều kiện cần\rq\rq, \lq\lq điều kiện đủ\rq\rq~để phát biểu lại định lí: \lq\lq Nếu tứ giác $ABCD$ là hình chữ nhật thì hai đường chéo bằng nhau\rq\rq.
\loigiai{
	Ta có thể phát biểu lại định lí đã cho như sau:\\
	\lq\lq Tứ giác $ABCD$ có hai đường chéo bằng nhau là điều kiện cần để nó là hình chữ nhật\rq\rq~hoặc \lq\lq Tứ giác $ABCD$ là hình chữ nhật là điều kiện đủ để hai đường chéo bằng nhau\rq\rq.
	}
\end{bt}
\begin{bt}%[0D1V1-4]
	Sử dụng thuật ngữ \lq\lq điều kiện cần và đủ\rq\rq, phát biểu lại các định lí sau:
	\begin{enumerate}[a)]
		\item Một phương trình bậc hai có hai nghiệm phân biệt khi và chỉ khi biệt thức của nó dương;
		\item Một hình bình hành là hình thoi thì nó có hai đường chéo vuông góc với nhau và ngược lại.
	\end{enumerate}
\loigiai{
	\begin{enumerate}[a)]
		\item Biệt thức của phương trình bậc hai dương là điều kiện cần và đủ để phương trình có hai nghiệm phân biệt.
		\item Hình bình hành có hai đường chéo vuông góc với nhau là điều kiện cần và đủ để hình bình hành là hình thoi.
	\end{enumerate}
	}
\end{bt}
\baitaptn
\TN
\setcounter{ex}{0}
\Opensolutionfile{ans}[ans/ans-c1b1-lc]
\begin{ex}%[Trích đề thi giữa học kì 1 - THPT MARIE CURIE - TPHCM - Năm học 2024-2025]%[0D1N1-1]
	Trong các câu sau, câu nào là mệnh đề?
	\choice
	{Hôm nay là thứ mấy?}
	{Hôm nay trời đẹp quá!}
	{\True Việt Nam là một quốc gia có bờ biển}
	{Hoa hồng đẹp nhất trong các loài hoa}
\loigiai{
	\lq\lq Việt Nam là một quốc gia có bờ biển\rq\rq\,là một mệnh đề.
	}
\end{ex}
\begin{ex}%[Trích đề giữa học kì 1 - THPT TÂN BÌNH - Tp HCM - Năm học 24-25]%[0D1H1-1]
	Trong các câu sau có bao nhiêu câu là mệnh đề?\\
	(1)$\colon$ Số $3$ là một số chẵn.\\
	(2)$\colon$ $2x+1=3$.\\
	(3)$\colon$ Các em hãy cố gắng làm bài thi cho tốt nhé!\\
	(4)$\colon$ $1<5\Rightarrow 8<6$.
	\choice
	{\True$2$}
	{$3$}
	{$1$}
	{$4$}
\loigiai{
	Các ý $(1)$ và $(4)$ là mệnh đề.}
\end{ex}
\begin{ex}%[0D1N1-3]
	Cho mệnh đề $A\colon$ \lq\lq Nghiệm của phương trình $x^2-5=0$ là số hữu tỉ\rq\rq. Mệnh đề phủ định của mệnh đề trên là
	\choice
	{\True \lq\lq Nghiệm của phương trình $x^2-5=0$ không là số hữu tỉ\rq\rq}
	{\lq\lq Nghiệm của phương trình $x^2-5=0$ không là số vô tỉ\rq\rq}
	{\lq\lq Phương trình $x^2-5=0$ vô nghiệm\rq\rq}
	{\lq\lq Nghiệm của phương trình $x^2-5=0$ không là số nguyên\rq\rq}
\loigiai{
	Mệnh đề phủ định của mệnh đề \lq\lq Nghiệm của phương trình $x^2-5=0$ là số hữu tỉ\rq\rq \ là \lq\lq Nghiệm của phương trình $x^2-5=0$ không là số hữu tỉ\rq\rq.
	}
\end{ex}
\begin{ex}%[Trích đề ôn tập giữa học kì 1 - noteên Lê Quý Đôn Ninh Thuận - Năm học 2024-2025]%[0D1N1-2]
	Cho mệnh đề \lq\lq Nếu hai tam giác bằng nhau thì diện tích chúng bằng nhau\rq\rq. Phát biểu nào sau đây đúng?
	\choice
	{Hai tam giác bằng nhau là điều kiện cần để diện tích chúng bằng nhau}
	{Hai tam giác bằng nhau là điều kiện cần và đủ để chúng có diện tích bằng nhau}
	{Hai tam giác có diện tích bằng nhau là điều kiện đủ để chúng bằng nhau}
	{\True Hai tam giác bằng nhau là điều kiện đủ để diện tích chúng bằng nhau}
\loigiai{Hai tam giác bằng nhau là điều kiện đủ để diện tích chúng bằng nhau.
	}
\end{ex}
\begin{ex}%[Trích đề ôn tập giữa học kì 1 - noteên Lê Quý Đôn Ninh Thuận - Năm học 2024-2025]%[0D1N1-5]
	Trong các mệnh đề dưới đây, mệnh đề nào đúng?
	\choice
	{$\forall x \in \mathbb{R},-x^2<0$}
	{$\forall x \in \mathbb{N}, x\,\vdots\,3$}
	{\True $\exists x \in \mathbb{R}, x>x^2$}
	{$\exists x \in \mathbb{R}, x^2-x+1 \leq 0$}
\loigiai{
	Với $x=\dfrac{1}{2}$, ta có $x^2=\dfrac{1}{4}$, khi đó $x> x^2$.}
\end{ex}
\begin{ex}%[0D1N1-5]
	Cho tứ giác $ABCD$. Xét mệnh đề \lq\lq Nếu tứ giác $ABCD$ là hình chữ nhật thì tứ giác $ABCD$ có hai đường chéo bằng nhau\rq\rq. Mệnh đề đảo của mệnh đề đó là
	\choice
	{\lq\lq Nếu tứ giác $ABCD$ là hình chữ nhật thì tứ giác $ABCD$ không có hai đường chéo bằng nhau\rq\rq}
	{\lq\lq Nếu tứ giác $ABCD$ không có hai đường chéo bằng nhau thì tứ giác $ABCD$ không là hình chữ nhật\rq\rq}
	{\lq\lq Nếu tứ giác $ABCD$ có hai đường chéo bằng nhau thì tứ giác $ABCD$ không là hình chữ nhật\rq\rq}
	{\True \lq\lq Nếu tứ giác $ABCD$ có hai đường chéo bằng nhau thì tứ giác $ABCD$ là hình chữ nhật\rq\rq}
\loigiai{
	Mệnh đề đảo cần tìm là \lq\lq Nếu tứ giác $ABCD$ có hai đường chéo bằng nhau thì tứ giác $ABCD$ là hình chữ nhật\rq\rq.
	}
\end{ex}
\begin{ex}%[0D1N1-5]
	Phủ định của mệnh đề \lq\lq $\exists x \in \mathbb{R}, x^2-x+1<0$\rq\rq \ là mệnh đề
	\choice
	{\True \lq\lq $\forall x \in \mathbb{R}, x^2-x+1 \geq 0$\rq\rq}
	{\lq\lq $\forall x \in \mathbb{R}, x^2-x+1<0$\rq\rq}
	{\lq\lq $\forall x \in \mathbb{R}, x^2-x+1>0$\rq\rq}
	{\lq\lq $\exists x \in \mathbb{R}, x^2-x+1 \geq 0$\rq\rq}
\loigiai{
	Phủ định của mệnh đề \lq\lq $\exists x \in \mathbb{R}, x^2-x+1<0$\rq\rq \ là mệnh đề \lq\lq $\forall x \in \mathbb{R}, x^2-x+1 \geq 0$\rq\rq.
	}
\end{ex}
\begin{ex}%[0D1N1-5]
	Phủ định của mệnh đề \lq\lq $\forall x \in \mathbb{R}, x^2 \geq 0$\rq\rq \ là mệnh đề
	\choice
	{\lq\lq $\exists x \in \mathbb{R}, x^2 \geq 0$\rq\rq}
	{\lq\lq $\exists x \in \mathbb{R}, x^2>0$\rq\rq}
	{\lq\lq $\exists x \in \mathbb{R}, x^2 \leq 0$\rq\rq}
	{\True \lq\lq $\exists x \in \mathbb{R}, x^2<0$\rq\rq}
\loigiai{
	Phủ định của mệnh đề \lq\lq $\forall x \in \mathbb{R}, x^2 \geq 0$\rq\rq \ là mệnh đề \lq\lq $\exists x \in \mathbb{R}, x^2<0$\rq\rq.
	}
\end{ex}
\begin{ex}%[0D1N1-4]
	Cho số tự nhiên $n$. Xét mệnh đề \lq\lq Nếu số tự nhiên $n$ chia hết cho $4$ thì $n$ chia hết cho $2$\rq\rq. Mệnh đề đảo của mệnh đề đó là
	\choice
	{\lq\lq Nếu số tự nhiên $n$ chia hết cho $2$ thì $n$ không chia hết cho $4$\rq\rq}
	{\lq\lq Nếu số tự nhiên $n$ chia hết cho $4$ thì $n$ không chia hết cho $2$\rq\rq}
	{\True \lq\lq Nếu số tự nhiên $n$ chia hết cho $2$ thì $n$ chia hết cho $4$\rq\rq}
	{\lq\lq Nếu số tự nhiên $n$ không chia hết cho $2$ thì $n$ không chia hết cho 4\rq\rq}
\loigiai{
	Mệnh đề đảo của mệnh đề \lq\lq Nếu số tự nhiên $n$ chia hết cho $4$ thì $n$ chia hết cho $2$\rq\rq \ là \lq\lq Nếu số tự nhiên $n$ chia hết cho $2$ thì $n$ chia hết cho $4$\rq\rq.
	}
\end{ex}
\begin{ex}%[Trích đề giữa học kì 1 - THPT TÂN BÌNH - Tp HCM - Năm học 24-25]%[0D1N1-1]
	Cho mệnh đề chứa biến \lq\lq$P(n) \colon 3 - n > 0$\rq\rq \, với $n$ là số tự nhiên. Mệnh đề nào sau đây đúng?
	\choice
	{\True $P(2)$}
	{$P(4)$}
	{$P(5)$}
	{$P(3)$}
\loigiai{
	Thay $n = 2$ vào \lq\lq$P(n) \colon 3 - n > 0$\rq\rq \, ta được $P(2)\colon 3 - 2 > 0 \Rightarrow$ $P(2)$ đúng.
	}
\end{ex}
\begin{ex}%[Trích đề giữa học kì 1 - THPT TÂN BÌNH - Tp HCM - Năm học 24-25]%[0D1H1-4]
	Biết rằng $P \Rightarrow Q$ là mệnh đề đúng. Mệnh đề nào sau đây đúng?
	\choice
	{\True $Q$ là điều kiện cần để có $P$}
	{$Q$ là điều kiện đủ để có $P$}
	{$P$ là điều kiện cần để có $Q$}
	{$Q$ là điều kiện cần và đủ để có $P$}
\loigiai{
	Mệnh đề $Q$ là điều kiện cần và đủ để có $P$ là phát biểu đúng nhất trong trường hợp $P \Rightarrow Q$ đúng.
	}
\end{ex}
\begin{ex}%[Trích đề học kì 1 - THPT Nguyễn Thị Minh Khai TpHCM Năm học 24-25]%[0D1H1-2]
	Xét các khẳng định sau
	\begin{enumerate}[1)]
		\item $2024$ chia hết cho $3$.
		\item $\pi$ là số vô tỉ.
		\item Phương trình $3x^2-4x+1=0$ có nghiệm nguyên.
		\item Nếu hai tam giác có diện tích bằng nhau thì hai tam giác đó bằng nhau.
	\end{enumerate}
	Trong các khẳng định trên, hỏi có tất cả bao nhiêu mệnh đề đúng?
	\choice
	{$1$}
	{$3$}
	{\True $2$}
	{$4$}
\loigiai{
	\begin{enumerate}
		\item $2024$ chia hết cho $3$ là khẳng định sai vì $2024:3=674$ dư $2$.
		\item $\pi$ là số vô tỉ là khẳng định đúng.
		\item Phương trình $3x^2-4x+1=0$ có nghiệm $x=1$ và $x=\dfrac{1}{3}$ nên phương trình $3x^2-4x+1=0$ có nghiệm nguyên là khẳng định đúng.
		\item Nếu hai tam giác có diện tích bằng nhau thì hai tam giác đó bằng nhau là khẳng định sai.
	\end{enumerate}
	Vậy có $2$ khẳng định đúng.
	}
\end{ex}
\Closesolutionfile{ans}
% \indapan{6}{ans/ans-c1b1-lc}
\TNTF
\setcounter{ex}{0}
\Opensolutionfile{ans}[ans/ans-c1b1-ds]
\begin{ex}%[0D1H1-5]
	Xét tính \textbf{đúng}, \textbf{sai} của mỗi mệnh đề sau
	\choiceTF[t]
	{$\forall x \in \mathbb{R}, x^2>0$}
	{\True $\exists a \in \mathbb{Q}, a>a^2$}
	{\True $\forall n \in \mathbb{Z}, n^2+n+2$ chia hết cho $2$}
	{$\forall n \in \mathbb{N}, n(n+1)(n+2)$ không chia hết cho $3$}
\loigiai{
	\begin{itemchoice}
		\itemch Ta chọn $x=0$ thì $x^2=0>0$ là \textbf{sai}.
		\itemch Ta chọn $a=\dfrac{1}{2} \in \mathbb{Q}$ thì $a^2=\dfrac{1}{4}$ nên $a>a^2$ (đúng).
		\itemch Thật vậy $\forall n \in \mathbb{Z}, n^2+n+2=n(n+1)+2$, trong đó $n(n+1)$ là tích của hai số nguyên liên tiếp nên chia hết cho $2$, vì vậy $n(n+1)+2$ cũng chia hết cho $2$.
		\itemch  Ta cho $n=1$ thì $n(n+1)(n+2)=1 \cdot 2 \cdot 3=6$ chia hết cho $3$.
	\end{itemchoice}
	}
\end{ex}
\begin{ex}%[0D1H1-5]
	Cho mệnh đề $P:$ \lq\lq$\forall x\in\mathbb{R},\, x^2>0$\rq\rq. Xét tính đúng sai của các khẳng định dưới đây.
	\choiceTF[t]
	{\lq\lq Với mọi $x\in\mathbb{R},\, x^2>2$\rq\rq\, là mệnh đề chứa biến}
	{$P$ là mệnh đề đúng}
	{$\overline{P}$ là mệnh đề sai}
	{\True Phủ định cua mệnh đề $P$ là $\overline{P}:$ \lq\lq $\exists x \in\mathbb{R}\colon x^2\leq 0$\rq\rq}
\loigiai{
	\begin{itemchoice}
		\itemch Mệnh đề \lq\lq Với mọi $x\in\mathbb{R},\, x^2>2$\rq\rq\, không là mệnh đề chứa biến.
		\itemch $P$ là mệnh đề sai vì với $x=0$ thì $0^2=0$.
		\itemch Vì $P$ sai nên $\overline{P}$ đúng.
		\itemch Phủ định của mệnh đề $P:$ \lq\lq$\forall x\in\mathbb{R},\, x^2>0$\rq\rq\, là mệnh đề $\overline{P}:$ \lq\lq $\exists\in\mathbb{R}\colon x^2\leq 0$\rq\rq.
	\end{itemchoice}
	}
\end{ex}
\begin{ex}%[0D1H1-2]
	Cho $P(n)=n^2-6n+10$ với $n$ là số tự nhiên.
	\choiceTF[t]
	{$P(1)$ chia hết cho $3$}
	{$P(2)$ là số lẻ}
	{$P(2n)>P(n)-1$ với $n=1$}
	{\True Tồn tại số tự nhiên $n$ thỏa mãn điều kiện $\dfrac{2\cdot P(n)-1}{n-3}$ là số nguyên}
\loigiai{
	\begin{itemchoice}
		\itemch Vì $P(1)=5$ không chia hết cho $3$.
		\itemch Vì $P(2)=2$ là số chẵn.
		\itemch Vì
		\begin{itemize}
			\item
			$P(2n)>P(n)-1\Leftrightarrow 4n^2-12n+10>n^2-6n+10-1\Leftrightarrow 3n^2-6n+1>0$.
			\item Khi $n=1\Rightarrow$ VT$=3\cdot 1^2-6\cdot 1+1=-2<0$.
		\end{itemize}
		\itemch Vì
		$\dfrac{2P(n)-1}{n-3}=\dfrac{2\left(n^2-6n+10\right)-1}{n-3}=\dfrac{2\left(n-3\right)^2+1}{n-3}=2\left(n-3\right)+\dfrac{1}{n-3}$.\\
		Suy ra $\dfrac{2P(n)-1}{n-3}\in\mathbb{Z}\Leftrightarrow\left(n-3\right) \in \{-1;1\} \Leftrightarrow n\in \{2;4\}$.
	\end{itemchoice}
	}
\end{ex}
\begin{ex}%[0D1H1-2]
	Với mỗi mệnh đề sau, em hãy chọn Đ (đúng) hoăc S (sai).
	\choiceTF[t]
	{Một tam giác cân thì hai góc đều bằng $60^{\circ}$}
	{\True Tích của hai số tự nhiên là một số lẻ khi và chỉ khi cả hai số đều là số lẻ}
	{1 là số nguyên tố lẻ nhỏ nhất}
	{\True Hình bình hành có hai đường chéo vuông góc với nhau là hình thoi}
\loigiai{
	\begin{itemchoice}
		\itemch Một tam giác cân thì hai góc bằng nhau chứ khống nhất thiết bằng $60^{\circ}$.
		\itemch Tích của hai số tự nhiên là một số lẻ khi và chỉ khi cả hai số đều là số lẻ.
		\itemch $1$ là không phải số nguyên tố.
		\itemch Hình bình hành có hai đường chéo vuông góc với nhau là hình thoi.
	\end{itemchoice}
	}
\end{ex}
\Closesolutionfile{ans}
% \indapan{3}{ans/ans-c1b1-ds}
% \newpage
\TNSA
\setcounter{ex}{0}
\Opensolutionfile{ans}[ans/ans-c1b1-kq]
%%%=============EX_1=============%%%
\begin{ex}%[0D1V1-5]
	Xác định số mệnh đề \textbf{sai} trong các mệnh đề sau?
	\begin{multicols}{2}
		\begin{itemize}
			\item \lq\lq$\exists x\in\mathbb{R}\colon x^2-3x+2=0$\rq\rq.
			\item \lq\lq$\forall x\in\mathbb{R}\colon x^2\geq 0$\rq\rq.
			\item \lq\lq$\exists n\in\mathbb{N}\colon n^2=n$\rq\rq.
			\item \lq\lq$\forall n\in\mathbb{N}\colon n<2n$\rq\rq.
		\end{itemize}
	\end{multicols}
	\par
	\shortans[oly]{1}
\loigiai{
	\begin{itemize}
		\item Ta có
		\[x^2-3x+2=0\Leftrightarrow \hoac{&x=2\\
		&x=1.
		}\]
		nên mệnh đề \lq\lq$\exists x\in\mathbb{R}\colon x^2-3x+2=0$\rq\rq~là mệnh đề đúng.
		\item $x^2\geq 0$ với mọi $x\in\mathbb{R}$ nên mệnh đề \lq\lq$\forall x\in\mathbb{R}\colon x^2\geq 0$\rq\rq~là mệnh đề đúng.
		\item Với $n=0$ thì $n^2=n$ nên mệnh đề \lq\lq$\exists n\in\mathbb{N}\colon n^2=n$\rq\rq~là mệnh đề đúng.
		\item Với $n=0$ thì $n=2n$ nên mệnh đề \lq\lq$\forall n\in\mathbb{N}\colon n<2n$\rq\rq~là mệnh đề sai.
	\end{itemize}
	}
\end{ex}
%%%=============================%%%
%%%=============EX_2=============%%%
\begin{ex}%[0D1V1-1]
	Với giá trị nào của $x$ thì \lq\lq$x\in\mathbb{N}, x^2-1=0$\rq\rq~là mệnh đề đúng?
	\par
	\shortans[oly]{1}
	\loigiai
	{Vì $x\in\mathbb{N}$ mà $x^2-1=0$ nên $x=1$.}
\end{ex}
%%%=============================%%%
%%%=============EX_3=============%%%
\begin{ex}%[0D1V1-5]
	Xác định số mệnh đề đúng trong các mệnh đề sau?
	\begin{multicols}{2}
		\begin{itemize}
			\item \lq\lq$\exists n\in \mathbb{N}$, $n^3-n$ không chia hết cho $3$\rq\rq.
			\item \lq\lq$\forall x\in \mathbb{R}$, $x<4\Leftrightarrow x^2<16$\rq\rq.
			\item \lq\lq$\exists k\in \mathbb{Z}$, $k^2+k+1$ là một số chẵn\rq\rq.
			\item \lq\lq$\forall x\in \mathbb{Z},\,\dfrac{2x^3-6x^2+x-3}{2x^2+1}\in \mathbb{Z}$\rq\rq.
		\end{itemize}
	\end{multicols}
	\par
	\shortans[oly]{1}
	\loigiai
	{\begin{itemize}
	\item Với $n\in \mathbb{N}$ thì $n^3-n=(n-1)n(n+1)$ là tích của ba số tự nhiên liên tiếp nên luôn chia hết cho $3$ nên mệnh đề \lq\lq$\exists n\in \mathbb{N}$, $n^3-n$ không chia hết cho $3$\rq\rq~là mệnh đề sai.
	\item Với $x=-5$ thì $x<4$ nhưng $x^2=25>16$ nên mệnh đề \lq\lq$\forall x\in \mathbb{R}$, $x<4\Leftrightarrow x^2<16$\rq\rq~là một mệnh đề sai.
	\item Với $k\in \mathbb{Z}$ thì $k^2+k=k(k+1)$ là tích của hai số nguyên liên tiếp nên $k^2+k$ là số chẵn, suy ra $k^2+k+1$ là số lẻ với mọi $k\in \mathbb{Z}$. Vậy mệnh đề \lq\lq$\exists k\in \mathbb{Z}$, $k^2+k+1$ là một số chẵn\rq\rq~là một mệnh đề sai.
	\item Ta có $\dfrac{2x^3-6x^2+x-3}{2x^2+1}=\dfrac{2x^2(x-3)+(x-3)}{2x^2+1}=\dfrac{(x-3)(2x^2+1)}{2x^2+1}=x-3$.\\
	Vậy $x\in \mathbb{Z}$ thì $x-3\in \mathbb{Z}$ hay $\dfrac{2x^3-6x^2+x-3}{2x^2+1}\in \mathbb{Z}$ là mệnh đề đúng.
	\end{itemize}	}
\end{ex}
%%%=============================%%%
%%%=============EX_4=============%%%
%Từ ngân hàng. Câu 2525. File gốc: E:\2. TAI LIEU TOAN THPT\NGUON LOC 10 - CAU TRUC MOI\NGUON LOC 10 -NEW\0-B1.tex
\begin{ex}%[0D1V1-2]
	Có bao nhiêu số nguyên $n$ để $P\left( n \right)\colon$ \lq\lq$2n^3+n^2+7n+1$ chia hết cho $2n-1$\rq\rq\, là mệnh đề đúng?
	\shortans{$4$}
\loigiai{
	Ta có: $2{ n^3}+{ n^2}+7n+1=\left( { n^2}+n+4 \right)\left( 2n-1 \right)+5$\\
	$2{ n^3}+{ n^2}+7n+1$ chia hết cho $2n-1$ $\Leftrightarrow $ $5$ chia hết cho $2n-1$ $\Leftrightarrow \hoac{
	& 2n-1=1 \\
	& 2n-1=-1 \\
	& 2n-1=5 \\
	& 2n-1=-5
	}\Leftrightarrow \hoac{
	& n=1 \\
	& n=0 \\
	& n=3 \\
	& n=-2.
	}$\\
	Vậy có $4$ giá trị nguyên của $n$.
	}
\end{ex}
\begin{ex}%[0D1V1-2]
	Cho mệnh đề $P:$ \lq\lq $\exists x\in \mathbb{Z}|x^2-4x-5< 0$\rq\rq. Có bao nhiêu giá trị của $x$ để $P$ là mệnh đề đúng?
	\shortans{5}
\loigiai{
	Ta có
	\allowdisplaybreaks
	\begin{eqnarray*}
		x^2-4x-5< 0 &\Leftrightarrow& \left(x+1\right)\cdot \left(x-5\right) < 0\\
		&\Leftrightarrow&\hoac{&\heva{&x+1< 0\\&x-5> 0}\\&\heva{&x+1> 0\\&x-5< 0}}\\
		&\Leftrightarrow&\hoac{&x\in \varnothing\\&x\in \left(-1;5\right)}\\
		&\Leftrightarrow&x\in (-1;5)\stackrel{x\in \mathbb{Z}}{\to}x\in \{0;1;2;3;4\}.
	\end{eqnarray*}
	Vậy có $5$ giá trị thỏa mãn yêu cầu.
	}
\end{ex}
\begin{ex}[Mức độ 3]%[Trần Bá Huy, dự án 24-25 Bài giảng 10,11]%[0D1V1-5]
	Cho tam giác $ABC$ vuông tại $A$, có hai cạnh góc vuông là $n$ và $2n+1$ với $n>0$. Biết rằng \lq\lq $\exists n\in\mathbb{N},\, \triangle ABC$ có diện tích bằng $264$\rq\rq\, là mệnh đề đúng. Tức là tồn tại $n_0\in\mathbb{N}$ sao cho \lq\lq $\triangle ABC$ có diện tích bằng $264$\rq\rq\, là mệnh đề đúng. Tìm $n_0$.
	\shortans{$16$}
\loigiai{
	Tam giác $ABC$ có diện tích bằng $264$ khi
	\[\dfrac{n(2n+1)}{2}=264\Leftrightarrow 2n^2+n-528=0\Leftrightarrow \hoac{&n=16\\ &n=-\dfrac{33}{2}.}\]
	Vậy $n_0=16$ thỏa mãn yêu cầu bài toán.
	}
\end{ex}
%%%=============================%%%
\Closesolutionfile{ans}
% \indapan{4}{ans/ans-c1b1-kq}
